\section{Podłączenie oraz sterowanie silnikami krokowymi}
Na samym początku warto wspomnieć o głównej zasadzie działania silnika elektrycznego. Posiada on zdolność do obracania się o konkretne kroki w odpowiedzi na impulsy elektryczne, co czyni go chętnie wybieranym sposobem wytwarzania napędu w precyzyjnych aplikacjach. Dzięki takiemu wariantowi sterowania, można dokładnie kontrolować swoją pozycję i prędkość. W zależności od sposobu zasilania uzwojeń statora, wirnik obraca się o stały kąt przy każdym impulsie w kierunku zdefioniowanym przez przeływający wewnątrz prąd. 
\newline
Można wyróżnić kilka dostępnych opcji wyprowadzeń silnka krokowego oraz powiązanym z tym podłączeniem ich do sterowników. Do wspomnianych sterowników silników krokowych jesteśmy w stanie podłączyć silniki z czterema, sześcioma, a nawet ośmioma wyprowadzeniami, czyli silniki bipolarne, unipolarne lub uniwersalne. Ma to wpływ na rezystancję, indukcyjność i prąd zwojeń, dzięki czemu kontrolujemy bardziej moment lub wpływamy na pobór mocy. Zaciski sterownika są oznaczone jako \textit{A},\textit{/A} dla jednej fazy i \textit{B},\textit{/B} dla drugiej, a kierunek obrotów silnika można określać, zamieniając miejscami przewody jednej fazy. Tą funkcję aktualnie spełnia pin \textit{Direction} widoczny na zewnątrz sterowników, który po podaniu odpowiednio stanu napięcia wyznacza kierunek obrotu. Warto też zwrócić uwagę na pozostałe wyprowadzenia przygotowane przez producentów konieczne do prawidłowej pracy silników, czyli kolejno \textit{Enable} oraz \textit{CLK}. Pierwszy służy do podawania zezwolenia na ruch silnika stanowiąc pewnego rodzaju ograniczenie, natomiast pin zegarowy umożliwiwa podawanie sygnału PWM z częstotliwością wpływającą na prędkość obrotów urządzenia.

